\documentclass{amsart}
% For dealing with references we use the comment environment
\usepackage{verbatim}
\newenvironment{reference}{\comment}{\endcomment}
%\newenvironment{reference}{}{}
\newenvironment{slogan}{\comment}{\endcomment}
\newenvironment{history}{\comment}{\endcomment}

% For commutative diagrams we use Xy-pic
\usepackage[all]{xy}

% We use 2cell for 2-commutative diagrams.
\xyoption{2cell}
\UseAllTwocells

% We use multicol for the list of chapters between chapters
\usepackage{multicol}

% This is generally recommended for better output
\usepackage{lmodern}
\usepackage[T1]{fontenc}

% For cross-file-references

% Package for hypertext links:
\usepackage{hyperref}

% For any local file, say "hello.tex" you want to link to please

% Theorem environments.
%
\theoremstyle{plain}
\newtheorem{theorem}[subsection]{Theorem}
\newtheorem{proposition}[subsection]{Proposition}
\newtheorem{lemma}[subsection]{Lemma}

\theoremstyle{definition}
\newtheorem{definition}[subsection]{Definition}
\newtheorem{example}[subsection]{Example}
\newtheorem{exercise}[subsection]{Exercise}
\newtheorem{situation}[subsection]{Situation}

\theoremstyle{remark}
\newtheorem{remark}[subsection]{Remark}
\newtheorem{remarks}[subsection]{Remarks}

\numberwithin{equation}{subsection}

% Macros
%
\def\lim{\mathop{\mathrm{lim}}\nolimits}
\def\colim{\mathop{\mathrm{colim}}\nolimits}
\def\Spec{\mathop{\mathrm{Spec}}}
\def\Hom{\mathop{\mathrm{Hom}}\nolimits}
\def\Ext{\mathop{\mathrm{Ext}}\nolimits}
\def\SheafHom{\mathop{\mathcal{H}\!\mathit{om}}\nolimits}
\def\SheafExt{\mathop{\mathcal{E}\!\mathit{xt}}\nolimits}
\def\Sch{\mathit{Sch}}
\def\Mor{\mathop{\mathrm{Mor}}\nolimits}
\def\Ob{\mathop{\mathrm{Ob}}\nolimits}
\def\Sh{\mathop{\mathit{Sh}}\nolimits}
\def\NL{\mathop{N\!L}\nolimits}
\def\CH{\mathop{\mathrm{CH}}\nolimits}
\def\proetale{{pro\text{-}\acute{e}tale}}
\def\etale{{\acute{e}tale}}
\def\QCoh{\mathit{QCoh}}
\def\Ker{\mathop{\mathrm{Ker}}}
\def\Im{\mathop{\mathrm{Im}}}
\def\Coker{\mathop{\mathrm{Coker}}}
\def\Coim{\mathop{\mathrm{Coim}}}

% Boxtimes
%
\DeclareMathSymbol{\boxtimes}{\mathbin}{AMSa}{"02}

%
% Macros for moduli stacks/spaces
%
\def\QCohstack{\mathcal{QC}\!\mathit{oh}}
\def\Cohstack{\mathcal{C}\!\mathit{oh}}
\def\Spacesstack{\mathcal{S}\!\mathit{paces}}
\def\Quotfunctor{\mathrm{Quot}}
\def\Hilbfunctor{\mathrm{Hilb}}
\def\Curvesstack{\mathcal{C}\!\mathit{urves}}
\def\Polarizedstack{\mathcal{P}\!\mathit{olarized}}
\def\Complexesstack{\mathcal{C}\!\mathit{omplexes}}
% \Pic is the operator that assigns to X its picard group, usage \Pic(X)
% \Picardstack_{X/B} denotes the Picard stack of X over B
% \Picardfunctor_{X/B} denotes the Picard functor of X over B
\def\Pic{\mathop{\mathrm{Pic}}\nolimits}
\def\Picardstack{\mathcal{P}\!\mathit{ic}}
\def\Picardfunctor{\mathrm{Pic}}
\def\Deformationcategory{\mathcal{D}\!\mathit{ef}}

\setlength{\emergencystretch}{3em}
\begin{document}
\title[Stacks Project, Tag 0D2M]{374 --- Fibrewise ampleness yields uniform local twisted-cohomology vanishing and section lifting}
\maketitle
\noindent Copyright (C) 2005--2025 Johan de Jong. Stacks Project authors. Source: \href{https://stacks.math.columbia.edu/tag/0D2M}{Tag 0D2M}.

\noindent Editorial difficulty note (not author text). Difficulty H2: One must obtain a single twist bound across every infinitesimal thickening of the fibre. The proof packages the layers as a finite graded module over the associated graded base algebra, then combines exact-sequence induction, formal functions and Nakayama to transfer vanishing and section surjectivity back to the local base..

\noindent Editorial provenance note (not author text). History: excerpt from revision \texttt{a0444\allowbreak 6e57e\allowbreak c1fbc\allowbreak 252a8\allowbreak 71afc\allowbreak ec775\allowbreak 2fb28\allowbreak 07b14}, coherent.tex, lines 5973--6088. Reference presentation was adapted; original-author context and core support blocks were retained or added. The mathematical wording is unchanged. Licensed under GNU FDL 1.2 or later; no invariant sections or cover texts. See ../licenses/COPYING. Context blocks reproduce author definitions and supporting results; they are not additional counted problems.

\begin{lemma}
\label{lemma-ample-on-fibre}
Let $f : X \to Y$ be a proper morphism of schemes with $Y$ Noetherian.
Let $\mathcal{L}$ be an invertible $\mathcal{O}_X$-module.
Let $\mathcal{F}$ be a coherent $\mathcal{O}_X$-module.
Let $y \in Y$ be a point such that $\mathcal{L}_y$ is ample on $X_y$.
Then there exists a $d_0$ such that for all $d \geq d_0$ we have
$$
R^pf_*(\mathcal{F} \otimes_{\mathcal{O}_X} \mathcal{L}^{\otimes d})_y = 0
\text{ for }p > 0
$$
and the map
$$
f_*(\mathcal{F} \otimes_{\mathcal{O}_X} \mathcal{L}^{\otimes d})_y
\longrightarrow
H^0(X_y, \mathcal{F}_y \otimes_{\mathcal{O}_{X_y}} \mathcal{L}_y^{\otimes d})
$$
is surjective.
\end{lemma}

\begin{proof}
Note that $\mathcal{O}_{Y, y}$ is a Noetherian local ring.
Consider the canonical morphism
$c : \Spec(\mathcal{O}_{Y, y}) \to Y$, see
Schemes, Equation (\href{https://stacks.math.columbia.edu/tag/02NA}{equation canonical morphism (Tag 02NA)}).
This is a flat morphism as it identifies local rings.
Denote momentarily $f' : X' \to \Spec(\mathcal{O}_{Y, y})$
the base change of $f$ to this local ring. We see that
$c^*R^pf_*\mathcal{F} = R^pf'_*\mathcal{F}'$ by
Lemma \href{https://stacks.math.columbia.edu/tag/02KH}{lemma flat base change cohomology (Tag 02KH)}.
Moreover, the fibres $X_y$ and $X'_y$ are identified.
Hence we may assume that $Y = \Spec(A)$ is the spectrum of
a Noetherian local ring $(A, \mathfrak m, \kappa)$ and $y \in Y$
corresponds to $\mathfrak m$. In this case
$R^pf_*(\mathcal{F} \otimes_{\mathcal{O}_X} \mathcal{L}^{\otimes d})_y =
H^p(X, \mathcal{F} \otimes_{\mathcal{O}_X} \mathcal{L}^{\otimes d})$
for all $p \geq 0$. Denote $f_y : X_y \to \Spec(\kappa)$ the projection.

\medskip\noindent
Let $B = \text{Gr}_\mathfrak m(A) =
\bigoplus_{n \geq 0} \mathfrak m^n/\mathfrak m^{n + 1}$.
Consider the sheaf $\mathcal{B} = f_y^*\widetilde{B}$
of quasi-coherent graded $\mathcal{O}_{X_y}$-algebras.
We will use notation as in Section \href{https://stacks.math.columbia.edu/tag/02O7}{section theorem formal functions (Tag 02O7)}
with $I$ replaced by $\mathfrak m$.
Since $X_y$ is the closed subscheme of $X$ cut out by
$\mathfrak m\mathcal{O}_X$ we may think of
$\mathfrak m^n\mathcal{F}/\mathfrak m^{n + 1}\mathcal{F}$
as a coherent $\mathcal{O}_{X_y}$-module, see
Lemma \href{https://stacks.math.columbia.edu/tag/087T}{lemma i star equivalence (Tag 087T)}. Then
$\bigoplus_{n \geq 0} \mathfrak m^n\mathcal{F}/\mathfrak m^{n + 1}\mathcal{F}$
is a quasi-coherent graded $\mathcal{B}$-module of finite type
because it is generated in degree zero over $\mathcal{B}$
and because the degree zero part is
$\mathcal{F}_y = \mathcal{F}/\mathfrak m \mathcal{F}$
which is a coherent $\mathcal{O}_{X_y}$-module.
Hence by Lemma \href{https://stacks.math.columbia.edu/tag/0897}{lemma graded finiteness (Tag 0897)} part (2)
we see that
$$
H^p(X_y, \mathfrak m^n \mathcal{F}/ \mathfrak m^{n + 1}\mathcal{F}
\otimes_{\mathcal{O}_{X_y}} \mathcal{L}_y^{\otimes d}) = 0
$$
for all $p > 0$, $d \geq d_0$, and $n \geq 0$. By
Lemma \href{https://stacks.math.columbia.edu/tag/089W}{lemma relative affine cohomology (Tag 089W)}
this is the same as the statement that
$
H^p(X, \mathfrak m^n \mathcal{F}/ \mathfrak m^{n + 1}\mathcal{F}
\otimes_{\mathcal{O}_X} \mathcal{L}^{\otimes d}) = 0
$
for all $p > 0$, $d \geq d_0$, and $n \geq 0$.

\medskip\noindent
Consider the short exact sequences
$$
0 \to \mathfrak m^n\mathcal{F}/\mathfrak m^{n + 1} \mathcal{F}
\to \mathcal{F}/\mathfrak m^{n + 1} \mathcal{F}
\to \mathcal{F}/\mathfrak m^n \mathcal{F} \to 0
$$
of coherent $\mathcal{O}_X$-modules. Tensoring with $\mathcal{L}^{\otimes d}$
is an exact functor and we obtain short exact sequences
$$
0 \to
\mathfrak m^n\mathcal{F}/\mathfrak m^{n + 1} \mathcal{F}
\otimes_{\mathcal{O}_X} \mathcal{L}^{\otimes d}
\to \mathcal{F}/\mathfrak m^{n + 1} \mathcal{F}
\otimes_{\mathcal{O}_X} \mathcal{L}^{\otimes d}
\to \mathcal{F}/\mathfrak m^n \mathcal{F}
\otimes_{\mathcal{O}_X} \mathcal{L}^{\otimes d} \to 0
$$
Using the long exact cohomology sequence and the vanishing above
we conclude (using induction) that
\begin{enumerate}
\item $H^p(X, \mathcal{F}/\mathfrak m^n \mathcal{F}
\otimes_{\mathcal{O}_X} \mathcal{L}^{\otimes d}) = 0$
for all $p > 0$, $d \geq d_0$, and $n \geq 0$, and
\item $H^0(X, \mathcal{F}/\mathfrak m^n \mathcal{F}
\otimes_{\mathcal{O}_X} \mathcal{L}^{\otimes d}) \to
H^0(X_y, \mathcal{F}_y \otimes_{\mathcal{O}_{X_y}} \mathcal{L}_y^{\otimes d})$
is surjective for all $d \geq d_0$ and $n \geq 1$.
\end{enumerate}
By the theorem on formal functions (Theorem \href{https://stacks.math.columbia.edu/tag/02OC}{theorem formal functions (Tag 02OC)})
we find that the $\mathfrak m$-adic completion of
$H^p(X, \mathcal{F} \otimes_{\mathcal{O}_X} \mathcal{L}^{\otimes d})$
is zero for all $d \geq d_0$ and $p > 0$.
Since $H^p(X, \mathcal{F} \otimes_{\mathcal{O}_X} \mathcal{L}^{\otimes d})$
is a finite $A$-module by
Lemma \href{https://stacks.math.columbia.edu/tag/02O6}{lemma proper over affine cohomology finite (Tag 02O6)}
it follows from Nakayama's lemma (Algebra, Lemma \href{https://stacks.math.columbia.edu/tag/00DV}{lemma NAK (Tag 00DV)})
that $H^p(X, \mathcal{F} \otimes_{\mathcal{O}_X} \mathcal{L}^{\otimes d})$
is zero for all $d \geq d_0$ and $p > 0$.
For $p = 0$ we deduce from
Lemma \href{https://stacks.math.columbia.edu/tag/02OB}{lemma ML cohomology powers ideal (Tag 02OB)} part (3)
that $H^0(X, \mathcal{F} \otimes_{\mathcal{O}_X} \mathcal{L}^{\otimes d}) \to
H^0(X_y, \mathcal{F}_y \otimes_{\mathcal{O}_{X_y}} \mathcal{L}_y^{\otimes d})$
is surjective, which gives the final statement of the lemma.
\end{proof}
\end{document}
